\section{Experiment 50: results}
\label{sec:exp50_results}

\subsection{The five readable coordinates preserve the measured CoG}

Without the 2-kpc fraction, the three quantile radii reproduce the CoG outside
roughly 5 kpc but leave the first measured aperture weakly constrained. Adding
$f_2$ lowers the median profile reconstruction RMS from $0.027\dex$ to
$0.00437\dex$. The adopted readable representation is $0.00100\dex$ worse than
the three-component PCA representation, whose median reconstruction RMS is
$0.00337\dex$, but both floors are more than an order of magnitude below the
halo-to-galaxy prediction error.

Adding $R_{90}$ lowers the readable reconstruction RMS to $0.00316\dex$, which
is $0.00021\dex$ better than PCA. It nevertheless changes held-out linear
profile CRPS from $0.06577\dex$ for the adopted five-coordinate representation
to $0.06601\dex$ for the six-coordinate version, making prediction worse by
$0.00023\dex$. The outermost additional coordinate is measurable from the CoG
but not sufficiently predictable from the available halo variables to improve
the reconstructed profile.

\begin{figure}[htbp]
  \centering
  \includegraphics[width=0.96\textwidth]{figures/exp50_representation.pdf}
  \caption{\ExpFifty{} representation study. The radial reconstruction errors
  and direct profile examples separate the information retained by each
  coordinate system from the later halo-prediction problem. The central-fraction
  coordinate is required to anchor the first measured aperture.}
  \label{fig:exp50_representation}
\end{figure}

\subsection{The readable halo-to-CoG map passes the PCA parity test}

The degree-2 readable model has held-out profile CRPS $0.06327\dex$. The
like-for-like degree-2 PCA reference has $0.06317\dex$. Thus the readable model
is worse by $0.00010\dex$, while the measured standard deviation among the five
PCA fold scores is $0.00024\dex$. It passes the gate in
\cref{eq:exp50_gate}.

The conditional-mean median profile RMS is $0.08560\dex$ for the readable model
and $0.08377\dex$ for PCA, so PCA is better by $0.00182\dex$ under this paired
mean statistic. The median largest fractional stellar-mass error beyond 5 kpc
is 24.64 per cent for the readable model and 24.25 per cent for PCA, making the
readable map worse by 0.39 percentage points. These small differences support
practical parity, not exact identity.

\begin{table}[htbp]
  \centering
  \caption{Held-out \ExpFifty{} prediction metrics. CRPS and RMS are in dex;
  lower values are better. Annular columns give RMS for CoG-derived stellar
  masses.}
  \label{tab:exp50_prediction}
  \resizebox{\textwidth}{!}{%
  \begin{tabular}{@{}lrrrrrr@{}}
    \toprule
    Model & Profile CRPS & Median profile RMS & $M(<10)$ & $M(10$--$30)$ & $M(30$--$50)$ & $M(50$--$100)$ \\
    \midrule
    Readable linear & 0.06577 & 0.08693 & 0.1168 & 0.1560 & 0.1717 & 0.1739 \\
    Readable degree-2 & 0.06327 & 0.08560 & 0.1115 & 0.1524 & 0.1670 & 0.1665 \\
    PCA linear & 0.06596 & 0.08631 & 0.1158 & 0.1555 & 0.1690 & 0.1697 \\
    PCA degree-2 & 0.06317 & 0.08377 & 0.1106 & 0.1513 & 0.1637 & 0.1667 \\
    \bottomrule
  \end{tabular}}
\end{table}

\begin{figure}[htbp]
  \centering
  \includegraphics[width=0.96\textwidth]{figures/exp50_predictive_skill.pdf}
  \caption{Held-out predictive comparison in \ExpFifty{}. The readable and PCA
  representations are evaluated on identical galaxy folds. The figure includes
  profile-wide scores and radial diagnostics; a low aggregate score is not used
  as a substitute for inspecting the reconstructed curves.}
  \label{fig:exp50_predictive_skill}
\end{figure}

\subsection{Total mass is much more predictable than detailed shape}

For the adopted degree-2 map, the held-out fraction of target-coordinate
variance explained is approximately 0.87 for total stellar mass, 0.31 for the
2-kpc fraction, 0.36 for $R_{20}$, 0.31 for the transformed $R_{20}$--$R_{50}$
gap, and 0.15 for the transformed $R_{50}$--$R_{80}$ gap. Therefore, the
selected halo variables tightly constrain overall stellar mass but only partly
constrain the radial mass distribution. A deterministic map would suppress
real profile diversity.

\begin{figure}[htbp]
  \centering
  \includegraphics[width=0.96\textwidth]{figures/exp50_parameter_predictions.pdf}
  \caption{Truth-versus-prediction and residual behavior for the five readable
  coordinates. The contrast between total-mass predictability and profile-shape
  predictability motivates treating the conversion as a conditional
  distribution.}
  \label{fig:exp50_parameters}
\end{figure}

\subsection{Assembly history and concentration both add paired information}

A final-mass-only linear model has median profile RMS $0.10890\dex$. Adding the
four \diffmah{} parameters lowers it to $0.09362\dex$, an improvement of
$0.01527\dex$. Shuffling the three MAH-shape variables at fixed final mass gives
$0.09472\dex$, which is $0.00109\dex$ worse than the correctly paired \diffmah{}
model. This is a modest but nonzero object-level assembly signal.

Adding correctly paired $z=0.4$ concentration lowers median profile RMS from
$0.09362\dex$ for \diffmah{} alone to $0.08693\dex$, an improvement of
$0.00669\dex$. Shuffling concentration within final-mass bins gives
$0.09338\dex$, nearly returning to the \diffmah{}-only reference. Concentration
therefore supplies information not encoded by the adopted smooth MAH
parameters.

\begin{figure}[htbp]
  \centering
  \includegraphics[width=0.96\textwidth]{figures/exp50_mass_binned_predictions.pdf}
  \caption{Measured and predicted CoGs with residuals in halo-mass bins at
  $z=0.4$. The conditional mean reproduces median trends but leaves broad
  object-to-object residuals, especially toward small radii.}
  \label{fig:exp50_mass_binned}
\end{figure}

\subsection{Correlated draws are required to reproduce the population}

The conditional mean is narrower than the TNG galaxy population. In the total
stellar mass--$R_{50}$ plane, its energy distance is 3.94 times the finite-sample
floor. One correlated draw per halo reduces the ratio to 1.35, making the draw
population 2.9 times closer to the floor under this ratio. In the central
$M_\star(<10\kpc)$ versus outer $M_\star(50$--$100\kpc)$ plane, the ratio falls
from 3.46 for the mean to 1.21 for a draw, again a 2.9-fold improvement. The
scatter model is thus necessary to represent the galaxy population.

The nominal 68 and 90 per cent pointwise intervals contain the measured CoGs
65.2 and 85.6 per cent of the time. The corresponding PCA coverages are 64.8
and 84.9 per cent. Both models are modestly under-covering, and 32 Monte Carlo
draws limit the precision of the tail estimates. The comparison does not reveal
a readable-coordinate-specific calibration failure.

\begin{figure}[htbp]
  \centering
  \includegraphics[width=0.96\textwidth]{figures/exp50_population_planes.pdf}
  \caption{Population-level validation in \ExpFifty{}. Conditional means
  collapse part of the intrinsic distribution, whereas correlated stochastic
  draws restore most of the observed width in mass--size and central--outer
  mass planes. Each population uses its own derived size.}
  \label{fig:exp50_population}
\end{figure}

\subsection{Symbolic regression finds useful corrections, but not a law}

The full linear map has profile CRPS $0.06577\dex$. Adding only nonlinear terms
that recur in at least four of five training folds lowers CRPS to
$0.06421\dex$, a 2.38 per cent improvement relative to linear. The dense
degree-2 model reaches $0.06327\dex$, a 3.80 per cent improvement; the sparse
terms therefore recover about 62 per cent of the available nonlinear gain.

The recurring standardized-variable terms are $c^2$ for total stellar mass,
$u_c\alpha_{\rm early}$ for the central fraction, and
$(\logten M_p)\alpha_{\rm early}$ for the middle radial gap. No correction
recurs in four folds for $R_{20}$ or the outer radial gap. Their coefficients
depend on the training-fold standardization, so these are empirical hypotheses,
not unit-independent physical equations.

\begin{figure}[htbp]
  \centering
  \includegraphics[width=0.96\textwidth]{figures/exp50_symbolic_regression.pdf}
  \caption{Nested-fold symbolic-regression result at $z=0.4$. Term recurrence
  is used as a stability filter, while the decoded held-out profile CRPS is the
  final criterion.}
  \label{fig:exp50_symbolic}
\end{figure}

\subsection{Concentration history adds a small independent signal}

On the 2366-galaxy overlap sample, the current-concentration degree-2 model has
profile CRPS $0.06318\dex$. Adding the five-epoch concentration history lowers
it to $0.06135\dex$, a 2.89 per cent improvement. Shuffling those histories
within final-mass bins gives $0.06325\dex$, which is 0.12 per cent worse than
the current-concentration reference. The real-history gain has the same sign
in all five held-out folds, whereas the shuffled history is worse in all five.

\begin{figure}[htbp]
  \centering
  \includegraphics[width=0.96\textwidth]{figures/exp50_concentration_history.pdf}
  \caption{Concentration-history extension in \ExpFifty{}. The real history is
  compared with current concentration and with histories permuted among halos
  of similar final mass. The gain is small but is not reproduced by the null.}
  \label{fig:exp50_concentration}
\end{figure}

\subsection{What \ExpFifty{} established}

At $z=0.4$, readable CoG coordinates can replace PCA coordinates without a
measurable loss at the precision of the five-fold comparison. The result is
probabilistic: only total stellar mass is tightly determined, and correlated
scatter is essential. Both MAH shape and concentration provide independent,
shuffle-controlled predictive information. The experiment does not establish
that one fixed relation applies at other epochs, that future halo history is
unnecessary, or that the sparse nonlinear terms are universal. Those questions
motivate \ExpFiftyOne{}.

\FloatBarrier

